\par\Needspace{6\baselineskip}\medskip
\begin{center}\hypertarget{native-182-article136}{}\pdfbookmark[2]{136.}{article-136}136.\end{center}
Zunächst sind die drei heliozentrischen Orte der Erde auf der Himmelskugel $A,A'\;A''$ (Abb. 4) mit
den drei entsprechenden geozentrischen Orten des Himmelskörpers $B,B',B''$ durch Großkreise zu
verbinden. Dann berechnet man sowohl die Lage dieser Großkreise in Bezug auf die Ekliptik
(sofern wir die Ekliptik als Grundebene wählen) als auch die Lage der Punkte $B,B',B''$ auf diesen
Kreisen. Es seien $\alpha,\alpha',\alpha''$ die drei geozentrischen Längen des Himmelskörpers, $\beta,\beta',\beta''$ seine Breiten
und $l,l',l''$ die heliozentrischen Längen der Erde, deren Breiten wir $=0$ setzen (Artikel 117,
72). Ferner seien $\gamma,\gamma',\gamma''$ die Neigungen der von $A,A',A''$ jeweils nach $B,B',B''$ gezogenen Großkreise
gegen die Ekliptik. Damit wir bei der Bestimmung dieser Neigungen nach einer festen Regel
verfahren, messen wir sie stets in Bezug auf den Teil der Ekliptik, der sich von den Punkten
$A,A',A''$ aus in der Reihenfolge der Tierkreiszeichen erstreckt. Ihre Größe wird also von $0$
bis $360^0$ gezählt, oder gleichbedeutend auf der Nordseite von $0$ bis $180^0$ und auf der
Südseite von $0$ bis $-180^0$. Die Bögen $AB,A'B',A''B''$, die wir stets zwischen $0$ und $180^0$
ansetzen können, bezeichnen wir mit $\delta,\delta',\delta''$. Zur Bestimmung von $\gamma,\delta$ haben wir damit die
Formeln
\[
\begin{aligned}
1)\quad&\operatorname{tang}\gamma=\frac{\operatorname{tang}\beta}{\sin(\alpha-l)}\\[.7ex]
2)\quad&\operatorname{tang}\delta=\frac{\operatorname{tang}(\alpha-l)}{\cos\gamma}\mathpunct{\raisebox{.45ex}{,}}
\end{aligned}
\]
Zur Kontrolle der Rechnung kann man bei Bedarf die folgenden Formeln hinzufügen
\[
\sin\delta=\frac{\sin\beta}{\sin\gamma}\mathpunct{\raisebox{.45ex}{,}}\qquad
\cos\delta=\cos\beta\cos(\alpha-l).
\]
Zur Bestimmung von $\gamma',\delta',\gamma'',\delta''$ gelten offenbar ganz entsprechende Formeln. Gälten $\beta=0$ und zugleich
entweder $\alpha-l=0$ oder $=180^0$, d.\,h. stünde der Himmelskörper in Opposition oder Konjunktion und
läge zugleich auf der Ekliptik, so wäre $\gamma$ unbestimmt. Wir setzen voraus, dass dieser Fall
bei keiner der drei Beobachtungen eintritt.

Wählt man als Grundebene anstelle der Ekliptik den Äquator, so benötigt man zur Bestimmung der
Lage der drei Großkreise in Bezug auf den Äquator außer den Neigungen auch die Rektaszensionen
der Schnittpunkte mit dem \hypertarget{native-183}{}Äquator. Außer den Abständen der Punkte $B,B',B''$ von diesen
Schnittpunkten sind auch die Abstände der Punkte $A,A',A''$ von ihnen zu berechnen. Da diese
Berechnungen von der in Artikel 110 behandelten Aufgabe abhängen, gehen wir hier auf die
Herleitung der Formeln nicht weiter ein.
